\clearpage
\begingroup
\fancyhead[L]{\small POSTHUMOUS PAPERS AND CORRESPONDENCE}
\renewcommand{\headrulewidth}{0pt}
\clearpage
\SourceLeaf{584}
\pdfbookmark[1]{Posthumous papers VIII. Perturbations of Pallas by Saturn}{pallas-nachlass-VIII}
\begin{center}
\large POSTHUMOUS PAPERS AND CORRESPONDENCE.
\par\bigskip
\Large [VIII. PERTURBATIONS OF PALLAS\par BY SATURN.]
\par\bigskip
\normalsize [1.]
\end{center}
\pdfbookmark[2]{[1]}{pallas-VIII-fragment1}
\begingroup
\small
\makeatletter\tagsleft@true\let\veqno\@@leqno\makeatother
\par\noindent\hspace*{\parindent}[The perturbations due to Saturn were calculated by \textsc{Nicolai} in essentially the same way
as those due to Jupiter. The perturbing forces were expanded by interpolation
in series with arguments \(M\) and \(M'\). The three quantities
\(T,V,W\) were set aside. The perturbing forces were resolved
into three rectangular components \(X,Y,Z\), of which the first lies along
the apsidal line of Pallas, the second lies perpendicular to it in the orbital plane of
Pallas, and the third is perpendicular to that plane. If \(x,y,z,x',y',z'\)
are the rectangular coordinates of the two planets referred to the axes just mentioned,
these components are
\[
\tag*{1)}
\begin{aligned}
X&=\mu\left(\frac{x'}{r'^3}-\frac{x'-x}{\rho^3}\right)\\[.6em]
Y&=\mu\left(\frac{y'}{r'^3}-\frac{y'-y}{\rho^3}\right)\\[.6em]
Z&=\mu\left(\frac{z'}{r'^3}-\frac{z'-z}{\rho^3}\right)\qquad(z=0)
\end{aligned}
\]
and their relations to the components used earlier are
\[
\begin{aligned}
\frac{ar}{\cos\varphi}\bigl(Y\cos(v-\tilde\omega)-X\sin(v-\tilde\omega)\bigr)&=V\\[.6em]
\frac{a}{\cos\varphi}\bigl(Y\sin(v-\tilde\omega)+X\cos(v-\tilde\omega)\bigr)&=T\\[.6em]
-\frac{ar}{\cos\varphi}Z&=W.
\end{aligned}
\]
\par\noindent\hspace*{\parindent}The perturbations of the elements are then determined from the equations
\SourceLeaf{585}
\[
\tag*{2)}
\begin{aligned}
\frac{\mathrm{d}i}{\mathrm{d}M}
&=-\frac{ar\cos(v-\AscendingNode)}{\cos\varphi}Z\\[.6em]
\frac{\mathrm{d}\AscendingNode}{\mathrm{d}M}
&=-\frac{ar\sin(v-\AscendingNode)}{\cos\varphi\sin i}Z\\[.6em]
\frac{\mathrm{d}\varphi}{\mathrm{d}M}
&=aa\sin(v-\tilde\omega)\mathord{.}\cos E\mathord{.}X
 -aa\bigl(1+\cos(v-\tilde\omega)\mathord{.}\cos E\bigr)Y\\[.6em]
\frac{\mathrm{d}\tilde\omega}{\mathrm{d}M}
&=\frac{aa}{e}\bigl(\cos\varphi+\sin(v-\tilde\omega)\mathord{.}\sin E\bigr)X
 -\frac{aa\cos(v-\tilde\omega)\sin E}{e}Y+A\\[.6em]
\frac{\mathrm{d}n}{\mathrm{d}M}
&=\frac{3a^3n\cos\varphi\cos E}{r}Y-\frac{3a^3n\sin E}{r}X\\[.6em]
\frac{\mathrm{d}\varepsilon}{\mathrm{d}M}
&=\frac{\mathrm{d}\tilde\omega}{\mathrm{d}M}
 -\frac{rr\left(\dfrac{\mathrm{d}\tilde\omega}{\mathrm{d}M}-A\right)}{aa\cos\varphi}
 -\frac{\sin E(2-ee-e\cos E)}{\cos\varphi}\frac{\mathrm{d}\varphi}{\mathrm{d}M}\\[.6em]
A&=-\frac{ar\,\mathrm{tg}\,\tfrac{1}{2}i\sin(v-\AscendingNode)}{\cos\varphi}Z,
\end{aligned}
\]
which are easily derived from Articles 8--12 of the Exposition d'une nouvelle méthode, pp. 450--455
of this volume. The notation has the meaning explained on pp. 473--474
of this volume.]
\par\endgroup
\clearpage
\endgroup
